\scanpage{20}
\mspage{8}
Il revient au $=$ de dire que $\sigma_0,\sigma_1,\sigma_\infty,\rho$ sont des générateurs, avec relations de définition
\[
(4)\qquad
\overbrace{(\rho^{-1}\sigma_\infty\sigma_0\deleted{\rho})}^{l_\infty}
\overbrace{(\rho^{-1}\sigma_1\sigma_\infty\deleted{\rho})}^{l_1}
\overbrace{(\rho^{-1}\sigma_0\sigma_1\deleted{\rho})}^{l_0}=1
\]
\[
(5)\qquad \sigma_0^2=\sigma_1^2=\sigma_\infty^2=1
\]
\[
(6)\qquad \rho^3=1.
\]
\margin{Il faut ajouter les relations $\rho\sigma_0\rho^{-1}=\sigma_1$ etc.
On a deux gén. $\sigma_0,\rho$ avec les relations
\[
\sigma_0^2=\rho^3=1,
\]
\deleted{$\unsure{[\rho^{-1}\sigma_0\rho\sigma_0\rho^{-1}]
[\underbrace{\rho\sigma_0\rho^{-1}}_{\sigma_1}
\underbrace{\rho^{-1}\sigma_0\rho}_{\sigma_\infty}\rho^{-1}]=1}$}
la relation (4) est conséquence des précédentes (\unsure{vient du fait} $\sigma_1=\rho\sigma_0\rho^{-1}$). Noter
\[
\begin{aligned}
l_0&=\deleted{\unsure{(\rho\sigma_0\rho^{-1})^2}}\rho^{-1}(\sigma_0\rho)^2\rho
=\boxed{\rho^{-1}\sigma_0\rho\sigma_0\rho^{-1}}\\
l_1&=(\deleted{\unsure{\rho}}\sigma_0\rho)^2\\
l_\infty&=(\rho^{-1}\sigma_0\deleted{\unsure{\rho}})^2
\end{aligned}
\]
\[
\begin{aligned}
\operatorname{int}(\sigma_0)l_1&=l_\infty\\
\operatorname{int}(\sigma_0)l_\infty&=l_1\\
\operatorname{int}(\sigma_0)(l_0)&=l_\infty^{-1}l_1^{-1}.
\end{aligned}
\]}
\margin{\uline{NB}
\[
\begin{aligned}
\sigma_0(l_0)&=l_\infty^{-1}l_1^{-1}\\
\sigma_0(l_1)&=l_\infty\\
\sigma_0(l_\infty)&=l_1\deleted{\unsure{l_\infty^{-1}}}
\end{aligned}
\]
et idem pour $\sigma_1,\sigma_\infty$.}
\uline{NB} Si on y fait \inserted{de plus} $\sigma_0,\sigma_1,\sigma_\infty=1$, on trouve comme relation de définition $\rho^3=1$, \deleted{donc} sans plus, ce qui prouve que $\rho$ n'est pas un mot \unsure{p.r.|p.v.} aux $\sigma_i$\ldots

Tout ceci est compatible avec les relations des débuts, \deleted{\unsure{si dans ces dernières on avait tout pris des primes}, et \unsure{pris les transposés} $l_i'=l_i^{-1}$, $\sigma_i'=\sigma_{\unsure{\dot\rho(i)}}$}
\deleted{i.e. $\sigma_i$} \deleted{\unsure{de notre groupe} $\sigma_i'$} si on y remplace les $l_i$ par les $l_i^{-1}$.
Ici on a, comme $\sigma_i=\sigma'_{\dot\rho^{-1}(i)}$
\[
(7)\quad\left\{\begin{aligned}
\operatorname{int}(\sigma_i)&:\text{ échange }l_j\text{ et }l_{\deleted{\unsure{j}}k}
\quad\inserted{(i,j,k)\text{ distincts}}\ \deleted{\text{et}}\\[-2pt]
&\deleted{\text{(et le }\unsure{\text{devenir}}\ l_{\dot\rho(i)}\ \text{s'en déduit par (1))}}\\
\operatorname{int}(\rho)l_i&=l_{\dot\rho(i)}.
\end{aligned}\right.
\]

On voudrait maintenant expliciter les $\varepsilon_i$ avec $\varepsilon_i^2=l_i$ $(i=0,1,\infty)$, $\varepsilon_i$ étant un des générateurs du centralisateur (ou normalisateur, c'est \unsure{idem}) des $L_i$ dans $\mathcal E^+$. (ou le centralisateur dans $\mathcal E$). Les calculs faits plus haut, avec les chgts. de notation indiqués, donnent
\deleted{\[
(8)\quad\left\{\begin{aligned}
\varepsilon_0&=l_0^{\unsure{-1}}\sigma_1\unsure{l_1^{-1}|l_0^{-1}}\\
\varepsilon_1&=l_\infty^{-1}\sigma_\infty l_1^{-1}\\
\varepsilon_\infty&=l_0^{-1}\sigma_0 l_\infty^{\unsure{-1}}
\end{aligned}\right.
\]}
\mspage{\unsure{19|10}}
